\section{Related Work}
\label{related_work}


\textbf{Uncertainty Estimation for Regression Tasks.} Deep learning regression models are increasingly required not only to provide accurate point estimates but also to quantify predictive uncertainty. A large body of research has focused on Bayesian neural networks, which place distributions over weights and approximate posterior inference using variational methods or Monte Carlo dropout, thereby producing predictive intervals \cite{gal2016dropout}. Another line of work employs ensembles of neural networks to capture both aleatoric and epistemic uncertainties, with randomized initialization or bootstrapped training providing diverse predictions \cite{lakshminarayanan2017simple}. More recently, post-hoc calibration techniques have been proposed, adapting classification-oriented approaches such as temperature scaling to regression settings, for instance by optimizing proper scoring rules or variance scaling factors \cite{kuleshov2018accurate}. Beyond probabilistic calibration, conformal prediction (CP) methods have gained attention due to their finite-sample coverage guarantees under minimal distributional assumptions. CP can be applied to regression to produce instance-dependent prediction intervals with guaranteed coverage, and has been extended to handle asymmetric intervals, distribution shift, and multi-target regression \cite{vovk2005conformal,cqr2019,Nizhar2025}. 

\textbf{Time Series Forecasting and Uncertainty Estimation.} Time series forecasting is a critical discipline in machine learning and statistics, focusing on predicting future values from a sequence of historical data points ordered by time. This field has wide-ranging applications, including financial market analysis, energy consumption forecasting, weather prediction, and medical prognosis. Traditional statistical methods, such as Autoregressive Integrated Moving Average (ARIMA) and Exponential Smoothing, have been foundational. However, their effectiveness is often limited by their assumption of linearity and their inability to capture complex, non-linear dependencies. More recently, deep learning models, employing  Transformers \cite{Yuqietal-2023-PatchTST,autoformer,reformer}, Multi-Layer Perceptrons (MLPs) \cite{wang2024timemixer, dlinear}, and Convolutional Neural Networks (CNNs) \cite{wu2023timesnet,wang2024timemixer}, were shown to be effective in modeling temporal dynamics and long-range dependencies \cite{wang2024deep,time_series_survey,xwang2024deep,wang2024timemixer}. The ability to quantify the uncertainty of a forecast, rather than providing just a single point estimate, is of paramount importance.  Uncertainty quantification provides a confidence interval for the prediction, which is crucial for risk management and informed decision-making. Some recent works have introduced uncertainty estimation to time series forecasting (see e.g. \cite{cini2025corel, wu2025eci}).
Given its wide-ranging applications, the importance of reporting uncertainty, and its challenging nature, time series forecasting serves as a highly suitable domain to evaluate the performance of MoGU. 

